\documentclass[aspectratio=169,11pt]{beamer}

% ============================================================
% presentation styling
% ============================================================

\usepackage[T1]{fontenc}
\usepackage[sfdefault]{carlito}
\renewcommand{\familydefault}{\sfdefault}

\usepackage{graphicx}
\usepackage{booktabs}
\usepackage{multirow}
\usepackage{amsmath,amssymb,mathtools,bm}
\usepackage{ragged2e}
\usepackage{adjustbox}
\usepackage{tikz}
\usetikzlibrary{positioning,arrows.meta,calc,fit,shapes.geometric}

% UCF colors
\definecolor{UCFGold}{HTML}{FFCA29}
\definecolor{UCFBlack}{HTML}{231F20}
\definecolor{SoftGray}{HTML}{666666}
\definecolor{LightGray}{HTML}{F3F3F3}

\setbeamercolor{normal text}{fg=black,bg=white}
\setbeamercolor{frametitle}{fg=black,bg=white}
\setbeamercolor{structure}{fg=UCFBlack}
\setbeamercolor{itemize item}{fg=UCFGold}
\setbeamercolor{itemize subitem}{fg=UCFGold}
\setbeamercolor{enumerate item}{fg=UCFBlack}

\setbeamertemplate{navigation symbols}{}
\setbeamertemplate{headline}{}
\setbeamertemplate{footline}{}

\setbeamerfont{frametitle}{size=\LARGE,series=\bfseries}
\setbeamerfont{normal text}{size=\normalsize}
\setbeamersize{text margin left=0.72cm,text margin right=0.72cm}

% --- UCF footer ---
\setbeamercolor{ucffooter}{bg=UCFGold,fg=UCFBlack}
\setbeamertemplate{footline}{%
  \leavevmode%
  \begin{beamercolorbox}[wd=\paperwidth,ht=0.27cm,dp=0cm]{ucffooter}%
    \raisebox{0pt}[0pt][0pt]{%
      \includegraphics[height=0.64cm]{assets/ucf_knight.png}%
    }%
    \hfill
    % frame number, so slides can be referred to by number
    \raisebox{0.06cm}{\scriptsize\insertframenumber}\hspace*{0.35cm}%
  \end{beamercolorbox}%
}

% Who presents the current section, shown at the right of every slide title.
% Usage: \presenter{Joshua} at the top of a section file.
\newcommand{\thepresenter}{}
\newcommand{\presenter}[1]{\renewcommand{\thepresenter}{#1}}

% Slide titles positioning
\setbeamertemplate{frametitle}{%
  \vspace*{0.18cm}
  \insertframetitle\hfill{\normalsize\mdseries\color{SoftGray}\thepresenter}\par
  \vspace*{0.06cm}
}

% Section divider helper.
% Usage: \sectiondivider{Motivation and Background}
\newcommand{\sectiondivider}[1]{%
  \begin{frame}
    \vfill
    \centering
    {\fontsize{31}{36}\selectfont\bfseries #1\par}
    \vfill
  \end{frame}
}

% A callout box for main takeaways.
% Usage: \begin{takeaway} ... \end{takeaway}
% (The colour is set before the box opens; setting it inside had no effect.)
\setbeamercolor{takeaway}{bg=LightGray,fg=black}
\newenvironment{takeaway}
  {\begin{beamercolorbox}[wd=\textwidth,sep=0.25cm,rounded=true]{takeaway}\justifying}
  {\end{beamercolorbox}}

% Source / citation line, right-aligned at the bottom of the slide (clear of the
% UCF knight). It is part of the text flow, so content never runs underneath it.
% Put it last in the frame (before any \note).
% Usage: \slidesource{Honeybee (Cha et al., CVPR 2024)}
\newcommand{\slidesource}[1]{%
  \par\vskip 0pt plus 1filll
  {\raggedleft\scriptsize\color{SoftGray}Source: #1\par}%
}

% One-line takeaway under a slide title: the point to get on first read.
% Usage: \slidepoint{MoE adds parameters without adding compute per token.}
\newcommand{\slidepoint}[1]{%
  {\large\color{SoftGray}#1\par}\vspace{0.25cm}%
}

% Box marking the team's own reading, as opposed to the paper's claims.
% Usage: \ourtake{...}
\setbeamercolor{ourtake}{bg=UCFGold!22,fg=black}
\newcommand{\ourtake}[1]{%
  \begin{beamercolorbox}[wd=\textwidth,sep=0.16cm,rounded=true]{ourtake}
    {\footnotesize\textbf{Our take:} #1\par}
  \end{beamercolorbox}%
}

% Small gray label above a block of content.
% Usage: \claimlabel{MM1's conclusions}
\newcommand{\claimlabel}[1]{{\footnotesize\color{SoftGray}#1\par}}

% A figure scaled to fit: never wider than the column, never taller than the given height.
% Usage: \slidefig[5cm]{Table01.png}   (default height 5.2cm)
\newcommand{\slidefig}[2][5.2cm]{%
  {\centering\includegraphics[width=\linewidth,height=#1,keepaspectratio]{#2}\par}%
}

% Placeholders for the parts still to fill in. They print in red so nothing is missed.
% Usage: \todo{final number of queries}    \figtodo[4cm]{training curve: loss and validation MRR}
%        \tbd for a table cell with no number yet
\newcommand{\todo}[1]{{\color{red}\textbf{[TODO: #1]}}}
\newcommand{\tbd}{{\color{red}\textbf{?}}}
\newcommand{\figtodo}[2][4.2cm]{%
  {\centering\fcolorbox{red}{white}{\parbox[c][#1][c]{0.92\linewidth}{\centering\color{red}\textbf{Figure TODO:} #2}}\par}%
}

% Figure path.
\graphicspath{{assets/}{figures/}}


% ============================================================
% HOW THIS DECK IS ORGANIZED
% ============================================================
% The talk is 8-9 minutes, then 3 minutes of questions. Everyone presents a part:
%
%   sections/a_problem.tex   Joshua (A)  problem, motivation, research gap, "same rhythm"  ~2:00
%   sections/b_data.tex      Junior (B)  overall pipeline, library, fake beatbox queries    ~2:15
%   sections/c_model.tex     Anika  (C)  encoder and contrastive training                   ~2:00
%   sections/d_results.tex   Dylan  (D)  evaluation, results, what's left                   ~2:15
%   sections/z_closing.tex   everyone    questions, references, backup slides
%
% Edit only your own file, so we don't overwrite each other on Overleaf.
% Anything still to fill in prints in red: \todo{...} for text, \figtodo{...} for figures.
% Speaker notes go in \note{...} inside a frame.
% Figures: put them in figures/ (PDF or PNG) and include with \slidefig[height]{name}.
% Helpers (\slidepoint, \slidefig, takeaway, \slidesource, \sectiondivider) are in preamble.tex.

\title{Hum a Beat, Find a Loop}
\subtitle{Retrieving drum loops by beatboxing them}
\author{Joshua Yu \quad Junior \quad Anika Nayak \quad Dylan Palmer}
\date{}

\begin{document}

% ============================================================
% TITLE SLIDE
% ============================================================
{
\setbeamertemplate{footline}{}

\begin{frame}[plain]
    \centering

    \vspace*{1.5cm}

    {\fontsize{31}{36}\selectfont\bfseries
    Hum a Beat, Find a Loop
    \par}

    \vspace{0.15cm}

    {\fontsize{20}{25}\selectfont
    Retrieving drum loops by beatboxing them
    \par}

    \vspace{0.65cm}

    {\fontsize{13}{16}\selectfont
    Joshua Yu \quad
    Junior \todo{surname} \quad
    Anika Nayak \quad
    Dylan Palmer
    \par}

    \vspace{0.25cm}

    {\fontsize{11}{14}\selectfont\color{gray}
    CAP 6614 Final Project
    \par}

    \vfill
\end{frame}
}

% ============================================================
% PRESENTATION SECTIONS
% ============================================================
% ============================================================
% A: Problem, motivation, research gap, "same rhythm"   (Joshua, ~2:00)
% ============================================================
\section{Problem and Motivation}
\presenter{Joshua}

\tikzset{
  abox/.style={draw=UCFBlack, rounded corners=3pt, fill=LightGray, minimum height=0.95cm,
               align=center, font=\small, inner sep=5pt},
  aplain/.style={abox, fill=white},
  ahot/.style={abox, fill=UCFGold!60},
  abad/.style={abox, fill=white, draw=SoftGray, text=SoftGray},
  aarr/.style={-{Stealth[length=1.8mm]}, thick, UCFBlack},
  anote/.style={font=\scriptsize, text=SoftGray, align=center},
  aloop/.style={draw=UCFBlack, fill=white, rounded corners=2pt, minimum width=3.1cm,
                minimum height=0.55cm, font=\footnotesize, anchor=west, inner sep=3pt},
}

% ------------------------------------------------------------
\begin{frame}{Find a drum loop by beatboxing it}
  \slidepoint{Sample libraries are searched with text tags. Tags say the genre, not where the kicks, snares and hats land.}

  \centering
  \begin{adjustbox}{max width=\textwidth}
  \begin{tikzpicture}[node distance=0.8cm]
    \node[aplain, minimum width=3.0cm] (q) {\textbf{``boom -- tss -- pf -- tss''}\\[-1pt]
      {\scriptsize 2 bars, beatboxed into a laptop mic}};
    \node[ahot, right=of q, minimum width=2.6cm] (sys) {Rhythm\\search};
    \node[aloop, right=1.0cm of sys, yshift=0.75cm, fill=UCFGold!25] (r1) {1.\ funk loop, kit A \hfill same rhythm};
    \node[aloop, below=0.12cm of r1.south west, anchor=north west] (r2) {2.\ funk loop, kit B \hfill same rhythm};
    \node[aloop, below=0.12cm of r2.south west, anchor=north west, draw=SoftGray, text=SoftGray] (r3) {3.\ soul loop \hfill close};
    \draw[aarr] (q) -- (sys);
    \draw[aarr] (sys.east) -- ++(0.45,0) |- (r1.west);
    \draw[aarr] (sys.east) -- ++(0.45,0) |- (r2.west);
    \draw[aarr] (sys.east) -- ++(0.45,0) |- (r3.west);
    \node[anote, below=0.25cm of q] {\textbf{Input:} an imperfect beatbox};
    \node[anote, below=0.25cm of r3] {\textbf{Output:} drum loops ranked\\by how well the rhythm matches};
  \end{tikzpicture}
  \end{adjustbox}

  \vspace{0.45cm}

  \begin{takeaway}
    People already share grooves by performing them. We want a search box that takes a beat instead of words:
    beatbox 2 bars, get back loops from the library that play the same rhythm, on any drum kit.
  \end{takeaway}

  \note{
    Producers dig through loop libraries using tags like ``funk'' or ``90 BPM''. Those tags describe the
    genre, not where the hits land. Musicians usually communicate a groove by performing it, by beatboxing
    or tapping. So our input is 2 bars of beatbox from a laptop mic, and the output is a ranked list of
    loops that play the same rhythm, whatever kit they're on.
  }
\end{frame}

% ------------------------------------------------------------
\begin{frame}{Same rhythm, completely different sound}
  \slidepoint{A loop and a beatbox of the same groove line up in time, but not in sound.}

  \slidefig[4.3cm]{timbre_gap.pdf}

  \vspace{0.15cm}
  \begin{columns}[T,onlytextwidth]
    \column{0.32\textwidth}
      \textbf{Timbre gap.} A mouth ``pf'' and a kick drum share timing, not spectrum.
    \column{0.32\textwidth}
      \textbf{Imperfect people.} Amateurs rush, drag, drop hi-hats and can't voice every hit.
    \column{0.32\textwidth}
      \textbf{Every voice differs.} So we test on people the model has never heard.
  \end{columns}

  \note{
    Top: one of our library loops. Bottom: a fake beatbox of the same 2 bars. The white lines are beats.
    The hits land in the same places, but the spectra look nothing alike: that's the timbre gap.
    [Optional demo: play the two clips, from results/listening/beatbox\_speed/audio/07\_drums.wav and
    07\_middle.wav.] On top of that, real people are sloppy, and everyone's kick sounds different.
  }
\end{frame}

% ------------------------------------------------------------
\begin{frame}{Research gap}
  \slidepoint{Query-by-beatboxing has been done by transcribing the beatbox first. Nobody has tested learned rhythm embeddings against that.}

  \centering
  \begin{adjustbox}{max width=\textwidth}
  \begin{tikzpicture}[node distance=0.55cm]
    \node[anchor=east, font=\small\bfseries, align=right] (lab1) at (0,0) {Classic\\{\scriptsize\mdseries Kapur et al.;}\\{\scriptsize\mdseries Nakano et al., 2004}};
    \node[abad, right=0.4cm of lab1] (a1) {Beatbox};
    \node[abad, right=of a1] (a2) {Label every\\sound};
    \node[abad, right=of a2] (a3) {Drum\\grid};
    \node[abad, right=of a3] (a4) {Match against\\library grids};
    \draw[aarr, SoftGray] (a1) -- (a2); \draw[aarr, SoftGray] (a2) -- (a3); \draw[aarr, SoftGray] (a3) -- (a4);
    \node[anote, right=0.25cm of a4, text=red!70!black, align=left] {one misheard\\sound breaks\\the match};

    \node[anchor=east, font=\small\bfseries, align=right] (lab2) at (0,-1.8) {Ours};
    \node[aplain, right=0.4cm of lab2] (b1) {Beatbox};
    \node[ahot, right=of b1, minimum width=3.6cm] (b2) {Learned\\rhythm embedding};
    \node[aplain, right=of b2] (b3) {Nearest loops,\\any kit};
    \draw[aarr] (b1) -- (b2); \draw[aarr] (b2) -- (b3);
    \node[anote, below=0.1cm of b2] {trained so beatbox and loop of\\the same rhythm land together};
  \end{tikzpicture}
  \end{adjustbox}

  \vspace{0.15cm}
  \raggedright\footnotesize
  \begin{itemize}
    \item \textbf{Close but different:} TRIA (2025) generates drums from beatboxes (no search); Ramires et al.\ (2018)
          adapt \emph{transcription} to each user; CLAP captures what audio means, not where hits land.
    \item \textbf{Our question:} do learned rhythm embeddings beat transcribe-then-match on people and drum kits never seen in training?
  \end{itemize}

  \note{
    The idea goes back to 2004: Kapur and colleagues, and Nakano and colleagues, transcribed each beatbox
    sound into a drum hit, then matched symbolically. If one sound is misheard, the match breaks. Recent work
    like TRIA uses beatboxes to generate drums rather than search, and per-user adaptation has only been
    done for transcription. So our gap: compare a learned rhythm embedding against transcribe-then-match,
    on held-out people and held-out kits.
  }
\end{frame}

% ------------------------------------------------------------
\begin{frame}{What counts as ``the same rhythm''?}
  \slidepoint{Each loop becomes a 3-voice grid. Same grid = same rhythm.}

  \slidefig[2.6cm]{rhythm_grid.pdf}

  \vspace{0.1cm}
  \begin{columns}[T,onlytextwidth]
    \column{0.55\textwidth}
      \small
      \begin{itemize}
        \item \textbf{3 voices:} kick, snare, hats (cymbals count as hats). Toms, crashes and ghost notes are left out.
        \item \textbf{12 steps per beat} hold both straight 16ths and triplets, so a swung beat is a different rhythm from a straight one.
      \end{itemize}
    \column{0.40\textwidth}
      \begin{takeaway}
        \centering
        {\Large\bfseries 8,626}\\
        distinct rhythms in 8,643 two-bar loops\\[2pt]
        {\footnotesize\color{SoftGray}Almost every loop is its own pattern, so each query has only a few right answers.}
      \end{takeaway}
  \end{columns}

  \note{
    This is the grid for the loop on the last slide. As the music lead, I wrote these rules: three voices,
    and twelve steps per beat, which is fine enough to tell straight 16ths from triplet swing. Two clips
    are the same rhythm only if the grids match exactly. Because both sides come from the same MIDI,
    the answer key is free. And 8,626 of the 8,643 loops are their own pattern, so this is close to
    finding one exact loop in the library.
  }
\end{frame}

% ============================================================
% B: Overall pipeline, library, fake beatbox queries   (Junior, ~2:15)
% ============================================================
% Scaffold: the pipeline figure and the facts below are filled in from our repo
% (spec v0.7, B's v0.5 handoff). Replace every red TODO, and cut whatever you
% don't need to fit the time.
\section{Method: Data}
\presenter{Junior}

\tikzset{
  pbox/.style={draw=UCFBlack, rounded corners=3pt, fill=LightGray, minimum height=1.0cm,
               align=center, font=\footnotesize, inner sep=4pt},
  pplain/.style={pbox, fill=white},
  phot/.style={pbox, fill=UCFGold!60},
  parr/.style={-{Stealth[length=1.8mm]}, thick, UCFBlack},
  pdash/.style={-{Stealth[length=1.8mm]}, thick, SoftGray, dashed},
  pnote/.style={font=\scriptsize, text=SoftGray, align=center},
  ptag/.style={circle, fill=UCFBlack, text=white, font=\tiny\bfseries, inner sep=1.2pt},
}
% Role tag on a box corner: \roletag{node}{A}
\newcommand{\roletag}[2]{\node[ptag, anchor=center] at (#1.north east) {#2};}

% ------------------------------------------------------------
\begin{frame}{Overall pipeline}
  \slidepoint{Library loops and fake queries are rendered from the same MIDI pattern, so every query's right answers are known exactly.}

  \centering
  \begin{adjustbox}{max width=\textwidth}
  \begin{tikzpicture}[node distance=0.55cm and 0.75cm]
    \node[pplain] (midi) {GMD drum\\pattern (MIDI)};
    \node[pbox, above right=0.1cm and 0.9cm of midi] (kits) {Render through\\7 drum kits};
    \node[pbox, below right=0.1cm and 0.9cm of midi] (bb) {AVP mouth sounds\\simplify + sloppiness};
    \node[pplain, right=of kits] (loop) {Library\\loop};
    \node[pplain, right=of bb] (fake) {Fake beatbox\\query};
    \node[pplain, below=0.45cm of fake] (real) {Real beatbox\\{\scriptsize team, test only}};
    \coordinate (mid) at ($(loop.east)!0.5!(fake.east)$);
    \node[phot, right=1.0cm of mid] (enc) {Shared\\rhythm encoder};
    \node[pbox, right=of enc] (emb) {128-d\\embedding};
    \node[pplain, right=of emb] (ret) {Top-$k$ loops\\by cosine};

    \draw[parr] (midi.east) -- ++(0.35,0) |- (kits.west);
    \draw[parr] (midi.east) -- ++(0.35,0) |- (bb.west);
    \draw[parr] (kits) -- (loop);
    \draw[parr] (bb) -- (fake);
    \draw[parr] (loop.east) -- (enc.west);
    \draw[parr] (fake.east) -- (enc.west);
    \draw[pdash] (real.east) -| (enc.south);
    \draw[parr] (enc) -- (emb);
    \draw[parr] (emb) -- (ret);

    \node[pnote, above=0.25cm of enc] {trained with SupCon:\\same pattern = positive pair};
    \node[pnote, below=0.12cm of midi] {grid label =\\pattern ID};

    \roletag{midi}{A}
    \roletag{bb}{A\,B}
    \roletag{kits}{B}
    \roletag{enc}{C}
    \roletag{ret}{D}
  \end{tikzpicture}
  \end{adjustbox}

  \vspace{0.2cm}
  {\footnotesize\color{SoftGray}
  \textbf{A} Joshua: musical rules (what counts as the same rhythm, how people beatbox) \quad
  \textbf{B} Junior: library and queries \quad
  \textbf{C} Anika: model \quad
  \textbf{D} Dylan: answer key and evaluation\par}

  \note{
    \todo{B: speaker notes.} Walk left to right: one MIDI pattern goes two ways, through drum kits to make
    library loops, and through real people's mouth sounds to make fake beatbox queries. Both go through one
    shared encoder; training pulls the same pattern together. At test time, real beatboxes go through
    the same encoder.
  }
\end{frame}

% ------------------------------------------------------------
\begin{frame}{The loop library}
  \slidepoint{Real drummers' MIDI, cut into 2-bar loops and played through drum kits, with separate kits for training and testing.}

  \begin{columns}[T,onlytextwidth]
    \column{0.5\textwidth}
      \small
      \begin{itemize}
        \item \textbf{Source:} Groove MIDI Dataset, 10 drummers, 388 performances.
        \item \textbf{8,643 two-bar loops} with GMD's own split: 6,904 train / 846 validation / 893 test.
        \item \textbf{7 free acoustic kits:} 5 for training (4 AVL kits, DRSKit),
              2 held out for test (Salamander, MuldjordKit).
        \item Recorded tempo, with a 0.25\,s lead-in.
        \item \todo{final rendered count (loops $\times$ kits) and hours of audio}
      \end{itemize}
    \column{0.46\textwidth}
      \figtodo[4.6cm]{one loop on two kits (spectrograms or waveforms), or a table of clips per split}
  \end{columns}

  \note{
    \todo{B: speaker notes.}
  }
\end{frame}

% ------------------------------------------------------------
\begin{frame}{Fake beatbox queries}
  \slidepoint{Each pattern is ``beatboxed'' with one real person's mouth sounds, simplified and made sloppy the way amateurs are.}

  \begin{columns}[T,onlytextwidth]
    \column{0.5\textwidth}
      \small
      \begin{itemize}
        \item \textbf{Voices:} AVP dataset, 28 amateur beatboxers, 9,776 usable kick, snare and hi-hat sounds.
              People split 20 train / 3 validation / 5 test.
        \item \textbf{Simplify like a person:} one sound at a time (kick, then snare, then hats);
              sounds never overlap; a speed limit (next slide).
        \item \textbf{Sloppiness:} steady rush or drag (up to $\pm$20\,ms), per-hit wobble (up to 40\,ms),
              dropped hits (hats up to 40\%), background noise (10--40\,dB SNR).
        \item \todo{number of fake queries per split}
      \end{itemize}
    \column{0.46\textwidth}
      \figtodo[4.6cm]{clean vs.\ sloppy beatbox of one pattern, with the drum loop (or play the audio)}
  \end{columns}

  \note{
    \todo{B: speaker notes.} Every value here comes from the music spec (config/music\_spec.yaml, v0.7).
  }
\end{frame}

% ------------------------------------------------------------
% Prepared by Joshua from the listening work. Present it here, or move it to
% Joshua's section if that flows better. Cut it if time is short.
\begin{frame}{Making the fake beatboxes sound human}
  \slidepoint{The first fake beatboxes were twice as fast as any real person. A speed limit, tuned by ear, fixed that.}

  \slidefig[3.4cm]{beatbox_speed.pdf}

  \vspace{0.1cm}
  \small
  \begin{itemize}
    \item Hits a drummer plays almost together landed on neighbouring grid steps and were voiced as separate sounds.
    \item \textbf{Speed limit:} a kick or snare needs 120\,ms from the sounds kept before it, a hat 250\,ms.
          Chosen by ear from five versions; 80\% of kicks and snares are still voiced.
    \item The answer key never changes: only what the fake beatboxer voices does.
  \end{itemize}

  \note{
    Compared over 1,500 random loops against 56 real AVP improvisations. Before: a median 5.9 sounds per
    second, and a quarter of gaps under 100 ms, which nobody can beatbox. With the limit: 3.4 per second,
    and the shortest gaps match real people's (about 130 ms).
  }
\end{frame}

% ============================================================
% C: Encoder and contrastive training   (Anika, ~2:00)
% ============================================================
% Scaffold: the facts below come from PR #1 (src/models/). Check them, replace
% every red TODO, and swap the practice numbers for v0 numbers when you have them.
\section{Method: Model}
\presenter{Anika}

\tikzset{
  cbox/.style={draw=UCFBlack, rounded corners=3pt, fill=LightGray, minimum height=1.0cm,
               align=center, font=\footnotesize, inner sep=4pt},
  cplain/.style={cbox, fill=white},
  chot/.style={cbox, fill=UCFGold!60},
  carr/.style={-{Stealth[length=1.8mm]}, thick, UCFBlack},
  cnote/.style={font=\scriptsize, text=SoftGray, align=center},
}

% ------------------------------------------------------------
\begin{frame}{One encoder for loops and beatboxes}
  \slidepoint{The same network embeds a drum loop and a beatbox, so ``same rhythm'' can become ``nearby vectors''.}

  \centering
  \begin{adjustbox}{max width=\textwidth}
  \begin{tikzpicture}[node distance=0.55cm]
    \node[cplain] (audio) {Audio\\{\scriptsize 22.05 kHz mono}};
    \node[cbox, right=of audio] (fe) {Front-end\\{\scriptsize R1: 3-band onsets}\\{\scriptsize R3: 64-bin log-mel}};
    \node[cbox, right=of fe] (conv) {Stride-2 conv\\{\scriptsize $\sim$20 ms tokens}};
    \node[chot, right=of conv] (tf) {4-layer\\Transformer\\{\scriptsize $d=128$}};
    \node[cbox, right=of tf] (pool) {Attention\\pooling};
    \node[cplain, right=of pool] (out) {128-d\\unit vector};
    \node[cplain, above=0.45cm of conv, font=\scriptsize] (kind) {kind token:\\library / beatbox / tap};
    \draw[carr] (audio) -- (fe); \draw[carr] (fe) -- (conv); \draw[carr] (conv) -- (tf);
    \draw[carr] (tf) -- (pool); \draw[carr] (pool) -- (out);
    \draw[carr] (kind.east) -| (tf.north);
    \node[cnote, below=0.15cm of tf] {weights shared by\\queries and loops};
  \end{tikzpicture}
  \end{adjustbox}

  \vspace{0.35cm}
  \raggedright
  \begin{itemize}
    \item About 550k parameters; trains on one GPU \todo{training time on v0 data}.
    \item Front-ends compared in experiment E2: R1 is timbre-light by design, R3 is a generic spectrogram.
          \todo{which front-end(s) the results use}
  \end{itemize}

  \note{
    \todo{C: speaker notes.}
  }
\end{frame}

% ------------------------------------------------------------
\begin{frame}{Contrastive training}
  \slidepoint{Clips of the same pattern are pulled together, whatever kit or voice plays them; other patterns are pushed apart.}

  \begin{columns}[T,onlytextwidth]
    \column{0.48\textwidth}
      \begin{itemize}
        \item \textbf{Supervised contrastive loss} (SupCon): the class is the pattern ID.
        \item \textbf{Batches} of $P$ patterns $\times$ $K$ views: different kits, different voices, different sloppiness.
              \todo{$P$, $K$, batch size}
        \item AdamW, cosine learning-rate decay, gradient clipping (some seeds collapsed without it).
        \item Validation MRR every epoch; the best checkpoint is kept.
      \end{itemize}
    \column{0.48\textwidth}
      \figtodo[4.8cm]{training curve: loss and validation MRR per epoch, on the v0 data}
  \end{columns}

  \note{
    \todo{C: speaker notes.}
  }
\end{frame}

% ------------------------------------------------------------
\begin{frame}{First check: the pipeline works}
  \slidepoint{On practice data, the encoder ranks the right loop far above chance.}

  \centering
  \begin{tabular}{lcccc}
    \toprule
    Front-end & MRR (95\% CI) & Hit@1 & Hit@5 & Hit@10 \\
    \midrule
    chance & 0.052 & & & \\
    R1: 3-band onsets & 0.373 (0.325--0.424) & 0.22 & 0.52 & 0.74 \\
    R3: log-mel & 0.793 (0.751--0.834) & 0.66 & 0.96 & 0.99 \\
    \bottomrule
  \end{tabular}

  \vspace{0.4cm}
  \raggedright\small
  \begin{itemize}
    \item Practice set: 100 test loops, 200 queries. The ``queries'' were the same loops on a second drum kit,
          not beatboxes, so R3 can partly match by sound. It shows the pipeline works; it doesn't settle R1 vs.\ R3.
    \item \todo{replace with, or add, the v0 numbers on fake beatbox queries}
  \end{itemize}

  \note{
    \todo{C: speaker notes.}
  }
\end{frame}

% ============================================================
% D: Evaluation, results, what's left   (Dylan, ~2:15)
% ============================================================
% Scaffold: the setup below follows docs/presentation_plan.md and README sections 7-9.
% Fill the red ? cells with Round 1 (Thu Nov 12) and Round 2 (Mon Nov 16) numbers,
% and cut any slide we have no result for. All numbers come from src/eval/evaluate.py.
\section{Experiments and Results}
\presenter{Dylan}

\tikzset{
  dbox/.style={draw=UCFBlack, rounded corners=3pt, fill=LightGray, minimum height=0.9cm,
               align=center, font=\footnotesize, inner sep=4pt},
  dplain/.style={dbox, fill=white},
  dhot/.style={dbox, fill=UCFGold!60},
  darr/.style={-{Stealth[length=1.8mm]}, thick, UCFBlack},
  dnote/.style={font=\scriptsize, text=SoftGray, align=center},
  drow/.style={draw=SoftGray, fill=white, rounded corners=2pt, minimum width=3.0cm,
               minimum height=0.42cm, font=\scriptsize, inner sep=2pt},
}

% ------------------------------------------------------------
\begin{frame}{How we evaluate}
  \slidepoint{Everything in the test is new to the model: new rhythms, new drum kits, new voices.}

  \begin{columns}[T,onlytextwidth]
    \column{0.52\textwidth}
      \small
      \begin{itemize}
        \item \textbf{Held out:} GMD test patterns (893 loops), 2 test kits, 5 AVP test people,
              and our own beatboxes.
        \item \textbf{Answer key:} same pattern ID. Each query has one right loop per test kit,
              among \todo{library size} clips.
        \item \textbf{Metrics:} mean reciprocal rank (MRR, main), Hit@1/5/10, mAP,
              bootstrap 95\% confidence intervals.
        \item \textbf{Compared with:} random chance and an onset baseline (no learning).
              \todo{CLAP if we get to it}
      \end{itemize}
    \column{0.44\textwidth}
      \centering
      \begin{tikzpicture}
        \node[dplain] (q) at (0,0) {Beatbox\\query};
        \node[drow] (r1) at (2.9,0.95) {1.\ loop 412};
        \node[drow] (r2) at (2.9,0.45) {2.\ loop 87};
        \node[drow, fill=UCFGold!60, draw=UCFBlack] (r3) at (2.9,-0.05) {3.\ loop 9 \hfill \textbf{right}};
        \node[drow] (r4) at (2.9,-0.55) {4.\ loop 230};
        \node[dnote] at (2.9,-0.95) {\dots};
        \draw[darr] (q.east) -- (1.35,0.2);
        \node[dnote, font=\small, text=UCFBlack] at (2.9,-1.55) {reciprocal rank $= 1/3$};
        \node[dnote] at (2.9,-2.05) {MRR = average over all queries};
      \end{tikzpicture}
  \end{columns}

  \note{
    \todo{D: speaker notes.}
  }
\end{frame}

% ------------------------------------------------------------
\begin{frame}{Results so far}
  \slidepoint{\todo{one sentence: the headline result}}

  \centering
  \begin{tabular}{lcccc}
    \toprule
    & \multicolumn{3}{c}{Fake beatbox (test people, test kits)} & Real beatbox (us) \\
    \cmidrule(lr){2-4}\cmidrule(lr){5-5}
    System & MRR (95\% CI) & Hit@1 & Hit@10 & MRR \\
    \midrule
    Random chance & \tbd & \tbd & \tbd & \tbd \\
    Onset baseline & \tbd & \tbd & \tbd & \tbd \\
    Our encoder & \tbd & \tbd & \tbd & \tbd \\
    \bottomrule
  \end{tabular}

  \vspace{0.35cm}
  \raggedright\small
  \begin{itemize}
    \item \todo{number of fake and real queries; which front-end ``our encoder'' is}
    \item \todo{what the gap between fake and real beatboxes tells us}
  \end{itemize}

  \note{
    \todo{D: speaker notes.} Round 1 numbers Thu Nov 12; real beatboxes Mon Nov 16.
  }
\end{frame}

% ------------------------------------------------------------
% Ablations. Keep whichever we have by Mon Nov 16 and cut the other.
\begin{frame}{What changes the result}
  \slidepoint{\todo{one sentence per finding}}

  \begin{columns}[T,onlytextwidth]
    \column{0.48\textwidth}
      \centering\textbf{Front-end}\par\smallskip
      \figtodo[4.4cm]{MRR for onset-envelope (R1) vs.\ log-mel (R3) input, fake and real beatboxes (C)}
    \column{0.48\textwidth}
      \centering\textbf{Messier queries}\par\smallskip
      \figtodo[4.4cm]{MRR as dropped hits and noise increase, encoder vs.\ onset baseline (stretch)}
  \end{columns}

  \note{
    \todo{D and C: speaker notes.}
  }
\end{frame}

% ------------------------------------------------------------
% Qualitative. Joshua (A) prepares the "which patterns fail" breakdown (a small E8)
% and can present this slide if that flows better.
\begin{frame}{What it gets right and wrong}
  \slidepoint{\todo{one sentence: the kind of rhythm it confuses}}

  \begin{columns}[T,onlytextwidth]
    \column{0.55\textwidth}
      \figtodo[4.6cm]{one real beatbox query and its top 5 loops as 3-voice grids, right answer highlighted (play it?)}
    \column{0.41\textwidth}
      \figtodo[4.6cm]{MRR split by pattern type: straight vs.\ swung, sparse vs.\ busy (A)}
  \end{columns}

  \note{
    \todo{speaker notes.}
  }
\end{frame}

% ------------------------------------------------------------
\begin{frame}{What's left}
  \slidepoint{The full plan answers three questions; the talk covers a first pass at the first one.}

  \begin{columns}[T,onlytextwidth]
    \column{0.32\textwidth}
      \begin{block}{Embedding vs.\ transcription?}
        \footnotesize
        \begin{itemize}
          \item The ``transcribe, then match'' baseline: the 2004 approach, rebuilt carefully
          \item CLAP baseline
          \item Tap queries, and real taps from TapTamDrum
        \end{itemize}
      \end{block}
    \column{0.32\textwidth}
      \begin{block}{Which input transfers to real people?}
        \footnotesize
        \begin{itemize}
          \item More front-ends, including a pretrained music encoder (MERT)
          \item Full recording sessions: about 50 beats each, with and without a click
          \item Fine-tune on 3 of us, test on the 4th
        \end{itemize}
      \end{block}
    \column{0.32\textwidth}
      \begin{block}{Does a quick enrollment help?}
        \footnotesize
        \begin{itemize}
          \item A new user records a few kicks, snares and hats first
          \item Adapt the encoder with LoRA, $k = 0, 1, 3, 5$ examples
        \end{itemize}
      \end{block}
  \end{columns}

  \vspace{0.25cm}
  {\footnotesize\color{SoftGray}\raggedright
  Also: electronic and lo-fi kits, other tempos, and tempo drift, swing drift and mic EQ in the fake beatboxes.\par}

  \note{
    \todo{D: speaker notes.}
  }
\end{frame}

% ============================================================
% Closing: questions, references, backup   (everyone)
% ============================================================
\presenter{}

% ------------------------------------------------------------
{
\setbeamertemplate{footline}{}
\begin{frame}[plain]
  \centering
  \vspace*{2.2cm}
  {\fontsize{31}{36}\selectfont\bfseries Questions?\par}
  \vspace{0.5cm}
  {\fontsize{13}{16}\selectfont
  Joshua Yu \quad Junior \todo{surname} \quad Anika Nayak \quad Dylan Palmer\par}
  \vspace{0.25cm}
  {\fontsize{11}{14}\selectfont\color{gray} Hum a Beat, Find a Loop \quad$\cdot$\quad CAP 6614\par}
  \vfill
\end{frame}
}

% ------------------------------------------------------------
% Add your section's sources here. Check each entry before the talk.
% If the list outgrows one slide, change this to \begin{frame}[allowframebreaks]{References}.
\begin{frame}{References}
  \scriptsize
  \begin{itemize}
    % --- Prior work (Joshua) ---
    \item A.~Kapur, M.~Benning, G.~Tzanetakis. \emph{Query-by-Beat-Boxing: Music Retrieval for the DJ.} ISMIR 2004.
    \item T.~Nakano, J.~Ogata, M.~Goto, Y.~Hiraga. \emph{A Drum Pattern Retrieval Method by Voice Percussion.} ISMIR 2004.
    \item O'Reilly et al. \emph{The Rhythm In Anything} (TRIA). ISMIR 2025.
    \item A.~Ramires, R.~Penha, M.~E.~P.~Davies. \emph{User Specific Adaptation in Automatic Transcription of Vocalised Percussion.} arXiv 2018.

    % --- Data (Junior) ---
    \item J.~Gillick, A.~Roberts, J.~Engel, D.~Eck, D.~Bamman. \emph{Learning to Groove with Inverse Sequence Transformations} (Groove MIDI Dataset). ICML 2019.
    \item A.~Delgado, S.~McDonald, N.~Xu, M.~Sandler. \emph{A New Dataset for Amateur Vocal Percussion Analysis} (AVP). Audio Mostly 2019.
    \item Drum kits: AVL Drumkits, DRSKit, Salamander Drumkit, MuldjordKit (free sample libraries). \todo{B: add each kit's license}

    % --- Model (Anika) ---
    \item P.~Khosla et al. \emph{Supervised Contrastive Learning.} NeurIPS 2020.

    % --- Evaluation (Dylan) ---
    \item Y.~Wu et al. \emph{Large-Scale Contrastive Language-Audio Pretraining with Feature Fusion and Keyword-to-Caption Augmentation} (LAION-CLAP). ICASSP 2023.

    % --- Remaining work ---
    \item B.~Haki et al. \emph{TapTamDrum: A Dataset for Dualized Drum Patterns.} ISMIR 2023.
    \item Y.~Li et al. \emph{MERT: Acoustic Music Understanding Model with Large-Scale Self-supervised Training.} ICLR 2024.
    \item E.~J.~Hu et al. \emph{LoRA: Low-Rank Adaptation of Large Language Models.} ICLR 2022.
  \end{itemize}

  \vspace{0.2cm}
  {\scriptsize\color{SoftGray} All diagrams and plots are our own.\par}
\end{frame}

% ------------------------------------------------------------
% Backup slides for questions go below. They come after References, so the
% talk itself ends on Questions.
\appendix


\end{document}
